\section{Cross problem synthesis and an evaluation agenda}
\label{sec:synthesis}

\subsection{The common object is a value with obligations}

Across the eight problems, a useful unit of reasoning is a value together with its obligations: the operations that produce it, the valid coordinates and numerical meaning, its consumers, its storage owner, and the events that permit reuse. A value is not exhausted by its payload bytes. A delayed memory request can require identity and return capacity before compute storage is bound. A reduced result can contain only a few logical elements while retaining a larger physical envelope. An asynchronously transmitted value can be locally complete yet still relevant to remote consumption or recovery.

A compiler can establish many obligations statically when control and access patterns are predictable. A kernel can encode them explicitly in thread roles, barriers, and buffers. Hardware can track them dynamically when readiness and resource availability vary. The architectural decision is where each obligation is most economically enforced while preserving a coherent contract between layers. This view avoids treating static scheduling, SIMT, dataflow, and superscalar issue as mutually exclusive answers to every part of an operator.

\subsection{Where the proposed mechanisms can help}

Fine-grained readiness is useful when consumers need independently valid subsets and when work released early has an available execution path. It provides little benefit if the next stage requires a global statistic, if every candidate waits on the same resource, or if the reduced coordination interval is smaller than the new bookkeeping cost. Partial readiness therefore needs both a semantic proof and a resource feasibility argument.

Late binding is useful when request latency occupies a substantial part of the early-reservation lifetime and when returns can be safely buffered or assigned capacity on arrival. It cannot eliminate the live payload needed by a consumer. Savings in final compute storage must be offset against return buffering, metadata, fragmentation, arbitration, and progress reserves. A system with a full return queue and no way to reclaim final storage can be worse than one with conservative early reservation.

Typed Tiles can carry the logical shape, representation, and axis of an operation, allowing software to avoid repeatedly rebuilding that meaning from individual lanes. This is an expressiveness and compilation benefit when the chosen operations are supported. It does not remove conversion, scale handling, packing, invalid tails, or numerical constraints. An exact-order reduction deliberately limits reassociation; an indexed atomic with returned old values deliberately limits aggregation. Such constraints must remain visible when judging potential throughput.

Independent scalar or event progress is useful when control work would otherwise serialize compute and communication. It must be paired with finite-resource allocation and fair shared-resource access. A second controller that waits on the same exhausted pool may add state without opening a drain path. Event offload that cannot recover or distinguish generations can trade a visible software protocol for a harder-to-debug hardware failure.

\subsection{A common experimental design}

A future matched evaluation should begin with a workload family rather than one favorable kernel. The dense family should vary independent output count, matrix tile shape, reduction length, dtype, layout, and the proportion of non-matrix work. The irregular family should vary index locality, duplicate rate, fan-out, sparse degree distribution, returned-value use, and supported atomic order. The distributed family should vary tile/message size, rank count, topology, imbalance, and service delays. These variables expose the conditions in the eight arguments without assuming that one mechanism is always active or useful.

The required outputs include end-to-end latency, steady-state throughput, and separate fill/drain contributions. Storage should be reported as an occupancy trace or live-interval distribution, not only a peak allocation constant. Movement should be counted at each interface, including conversion and staging. Control accounting should include dynamic instructions, events, descriptors, tag and queue state, selection/arbitration work, and time blocked for capacity. Physical comparisons require matched implementation technology and clock/area/power methodology; source-level counts cannot substitute for that evidence.

The numerical baseline must be explicit. Compare exact-order folds with exact-order folds, or state that reassociation is permitted and quantify the resulting error/reproducibility requirement. Distinguish FP32 storage from FP32 arithmetic and accumulation. The memory baseline must likewise be explicit: atomics, acquire/release scopes, result publication, safe reuse, and error recovery must satisfy equivalent obligations. A kernel that omits required synchronization is not a faster implementation of the same contract.

Ablations should disable or constrain one proposed mechanism at a time while keeping the rest of the workload and machine budget meaningful. Examples include whole-Tile versus partial readiness, early versus delayed binding with the same total storage accounting, fixed versus dynamic issue within the same finite window, and shared versus independent completion service. Enlarging every queue to conceal a liveness flaw is not a valid ablation. Each configuration must first satisfy safety and the progress property it claims.

\subsection{Correctness and progress before performance conclusions}

A useful validation hierarchy starts with executable instruction semantics and boundary cases, then compiler lowering and complete operator equivalence, then finite-capacity protocol testing, and finally performance/physical implementation. It includes shape and mask legality, type/encoding errors, collisions and old values, deterministic publication, buffer-generation reuse, exceptional numeric status, and memory-order litmus tests. Progress validation must include capacity saturation, delayed peers, simultaneous returns, fairness, cancellation, and the release path itself.

A passing finite test set is evidence, not a general proof. Where a design makes a broad claim such as deadlock freedom, it needs explicit invariants and environmental assumptions, with model checking or another defensible proof method for the relevant state space. Starvation freedom additionally needs an arbitration argument. A bounded latency claim requires bounded service and arrival assumptions, not just the absence of deadlock. These distinctions should survive into performance reports rather than being hidden in a single ``works'' label.

This article has not executed those model tests, compiler tests, benchmarks, or synthesis flows. It provides the contracts and questions that make such work interpretable. Approximate utilization or improvement percentages in prior advocacy material are excluded because the available provenance does not establish a matched baseline or a reproducible measurement. No novelty claim, architecture ranking, or whole-core PPA conclusion follows from the present evidence.

\subsection{Conclusion}

A Superscalar NPU is a plausible way to assign selected dependency, readiness, and storage decisions to bounded hardware. Its strongest argument is specific: a value-oriented execution model may reduce exposed coordination and conservative waiting in mixed, irregular, or variable-latency pipelines while preserving scalar programmability. The competing designs already solve many of these problems through optimized compilation, thread ownership, explicit roles, and asynchronous engines. The question is which placement of responsibility is most effective under a given workload and resource budget.

The eight problems make that decision concrete. Coordination must identify useful granularity; latency hiding must account for all live storage; representation must preserve numeric and layout meaning; irregularity must respect collision and returned-value contracts; reduction must respect physical ownership and permitted order; events must have a progressing agent; faults must have a defined state and visibility boundary; and every admitted operation must have a finite-resource completion path. A successful design closes all eight obligations together. Saving one wait, opcode, or software barrier is valuable only if the complete system remains correct, live, and less costly for the work it must perform.
